\documentclass[12pt,a4paper]{article}

% --- PACOTES BÁSICOS E TIPOGRAFIA ---
\usepackage[utf8]{inputenc}
\usepackage[T1]{fontenc}
\usepackage[brazil]{babel}
\usepackage{helvet}               % Fonte padrão Arial / Helvetica
\renewcommand{\familydefault}{\sfdefault}

% --- MARGENS (ABNT: Superior/Esquerda 3cm, Inferior/Direita 2cm) ---
\usepackage{geometry}
\geometry{
    a4paper,
    top=3cm,
    left=3cm,
    bottom=2cm,
    right=2cm
}

% --- ESPAÇAMENTOS ---
\usepackage{setspace}
\setstretch{1.5}                  % Espaçamento 1,5 entre linhas
\setlength{\parindent}{1.25cm}    % Recuo padrão de parágrafo
\setlength{\parskip}{0pt}         % Sem espaçamento extra entre parágrafos

% --- PAGINAÇÃO (ABNT: Superior Direito, tamanho 10pt) ---
\usepackage{fancyhdr}
\pagestyle{fancy}
\fancyhf{}
\renewcommand{\headrulewidth}{0pt}
\rhead{\fontsize{10}{12}\selectfont\thepage}

% --- PACOTES AUXILIARES PARA TABELAS E FORMATAÇÃO ---
\usepackage{tabularx}
\usepackage{array}
\usepackage{titlesec}
\usepackage{enumitem}

% Formatação dos Títulos de Seção (Caixa Alta e Negrito)
\titleformat{\section}
  {\normalfont\bfseries\normalsize\MakeUppercase}{\thesection}{1em}{}
\titleformat{\subsection}
  {\normalfont\bfseries\normalsize}{\thesubsection}{1em}{}
\titleformat{\subsubsection}
  {\normalfont\bfseries\normalsize}{\thesubsubsection}{1em}{}

% ==============================================================================
% DEFINIÇÃO DE VARIÁVEIS DO DOCUMENTO
% ==============================================================================
\newcommand{\varInstituicao}{CENTRO DE EDUCAÇÃO TECNOLÓGICA DO AMAZONAS -- CETAM}
\newcommand{\varCurso}{CURSO TÉCNICO EM DESENVOLVIMENTO DE SISTEMAS}
\newcommand{\varTitulo}{SISFEIRA: SISTEMA WEB DE GESTÃO DE PEDIDOS PARA FEIRA DE PRODUTORES LOCAIS}
\newcommand{\varEstudanteUm}{Paulo Monteiro}
\newcommand{\varEstudanteDois}{Luis Fernando}
\newcommand{\varEstudanteTres}{José Luiz}
\newcommand{\varMunicipio}{Manaus, AM}
\newcommand{\varAno}{2026}
\newcommand{\varProfessor}{Priscila Gonçalves}

% ==============================================================================
% INÍCIO DO DOCUMENTO
% ==============================================================================
\begin{document}

% ------------------------------------------------------------------------------
% 1. CAPA
% ------------------------------------------------------------------------------
\thispagestyle{empty}
\begin{center}
    {\bfseries\MakeUppercase{\varInstituicao}}\\[0.2cm]
    {\bfseries\MakeUppercase{\varCurso}}
    
    \vspace*{\fill}
    
    {\bfseries\varTitulo}
    
    \vspace{1.5cm}
    
    {\bfseries\varEstudanteUm}\\[0.1cm]
    {\bfseries\varEstudanteDois}\\[0.1cm]
    {\bfseries\varEstudanteTres}
    
    \vspace*{\fill}
    
    {\bfseries\varMunicipio}\\[0.1cm]
    {\bfseries\varAno}
\end{center}
\newpage

% ------------------------------------------------------------------------------
% FOLHA DE AVALIAÇÃO FINAL
% ------------------------------------------------------------------------------
\thispagestyle{empty}
\begin{center}
    {\bfseries\MakeUppercase{\varInstituicao}}\\[0.2cm]
    {\bfseries\MakeUppercase{\varCurso}}
    
    {\bfseries Avaliação Final do Projeto Técnico}
\end{center}

\vspace{0.8cm}
\noindent A avaliação do Projeto intitulado \textbf{\varTitulo}, desenvolvido pelos estudantes: \textbf{\varEstudanteUm}, \textbf{\varEstudanteDois} e \textbf{\varEstudanteTres}, é composta por duas partes:

\vspace{0.8cm}
\noindent \textbf{*Parte 1} -- Avaliação do acompanhamento sistemático realizado no decorrer do desenvolvimento do projeto, obtendo a nota: (\hspace{1.5cm})\underline{\hspace{3cm}}

\vspace{0.5cm}
\noindent \textbf{*Parte 2} -- Entrega do trabalho impresso e socialização da temática com a turma: (\hspace{1.5cm})\underline{\hspace{3cm}}

\vspace{1cm}
\begin{center}
    \textbf{Nota final nesta Etapa do Componente (Projeto): (\hspace{1.5cm})\underline{\hspace{3cm}}}
\end{center}

\vspace{1cm}
\noindent \textbf{Considerações do Professor/Instrutor:}
\vspace{0.3cm}

\noindent\rule{\textwidth}{0.4pt}\\[0.6cm]
\rule{\textwidth}{0.4pt}\\[0.6cm]
\rule{\textwidth}{0.4pt}\\[0.6cm]
\rule{\textwidth}{0.4pt}

\vspace{1.5cm}
\begin{center}
    \underline{\hspace{9cm}}\\
    \varProfessor\\
    Professor/Instrutor do Projeto (Avaliador)
    
    \vspace*{\fill}
    
    \varMunicipio\\
    \varAno
\end{center}
\newpage

% ------------------------------------------------------------------------------
% SUMÁRIO
% ------------------------------------------------------------------------------
\thispagestyle{empty}
\begin{center}
    {\bfseries\MakeUppercase{Sumário}}
\end{center}
\vspace{0.5cm}
\begin{spacing}{1.0}
\tableofcontents
\end{spacing}
\newpage

% Ativa numeração de páginas a partir do elemento textual
\setcounter{page}{4}

% ------------------------------------------------------------------------------
% 1. APRESENTAÇÃO
% ------------------------------------------------------------------------------
\section{Apresentação}

O projeto SISFEIRA consiste no desenvolvimento de uma aplicação web voltada à gestão de pedidos antecipados, catálogo de produtos e organização de retiradas para feiras livres de produtores locais. A proposta surge da necessidade de aproximar os agricultores familiares e pequenos produtores dos seus consumidores finais por meio de um canal digital acessível, direto e descomplicado.

Ao permitir que os clientes visualizem a oferta semanal e realizem seus pedidos previamente à realização física da feira, o sistema possibilita que os produtores colham e transportem apenas a quantidade demandada, mitigando perdas financeiras e desperdício de alimentos perecíveis.

A solução técnica prioriza a simplicidade operacional, o controle do código e a eficiência de execução, adotando tecnologias web abertas e de ampla portabilidade (HTML5, CSS3, JavaScript puro, Node.js, Express e PostgreSQL). O sistema é arquitetado em modelo monolítico modular com entrega estática \textit{Zero-Build} executada diretamente sobre o sistema operacional Debian Linux. Este documento apresenta a caracterização detalhada da proposta, a especificação formal de requisitos, as regras de negócio, a modelagem conceitual e de dados, bem como a metodologia e o cronograma de execução estabelecidos para a entrega do Produto Mínimo Viável (MVP) em um ciclo de 15 dias.

\newpage

% ------------------------------------------------------------------------------
% 2. CARACTERIZAÇÃO DO PROBLEMA E JUSTIFICATIVA TÉCNICA
% ------------------------------------------------------------------------------
\section{Caracterização do Problema e Justificativa Técnica}

As feiras livres representam um dos principais pilares de comercialização da agricultura familiar no Estado do Amazonas, viabilizando o sustento de centenas de famílias produtoras e garantindo o abastecimento de produtos frescos para os centros urbanos. No entanto, o modelo operacional tradicional dessas feiras enfrenta gargalos críticos decorrentes da inexistência de instrumentos digitais acessíveis para previsão de demanda e divulgação direta da colheita.

Historicamente, os produtores locais transportam suas mercadorias até as bancas com base em estimativas empíricas. Essa ausência de previsibilidade gera dois problemas recorrentes: o desabastecimento precoce de itens com alta procura ou o excedente de produtos perecíveis não comercializados que, diante das limitações logísticas e de refrigeração regional, acabam descartados como lixo orgânico ao término da jornada. Paralelamente, os consumidores enfrentam dificuldades para conhecer com antecedência a disponibilidade de itens sazonais ou reservar mercadorias diretamente com os feirantes de sua preferência.

A justificativa técnica do SISFEIRA fundamenta-se no emprego pragmático de software livre para viabilizar um canal de pedidos antecipados que otimize a cadeia de valor da feira:

\begin{itemize}[leftmargin=1.5cm]
    \item \textbf{Impacto Econômico:} Redução de prejuízos operacionais dos agricultores familiares, permitindo uma colheita dimensionada estritamente conforme o volume de pedidos confirmados.
    \item \textbf{Impacto Social e Digital:} Inclusão tecnológica dos feirantes por meio de interfaces leves e responsivas, adaptadas para uso em smartphones (\textit{Mobile-First}), sem a barreira de instalação de aplicativos complexos.
    \item \textbf{Impacto Ambiental:} Mitigação direta do desperdício de alimentos perecíveis e racionalização do frete e transporte das cargas das comunidades rurais para a área urbana.
\end{itemize}

Dessa forma, o projeto se posiciona como uma resposta prática, econômica e funcional, estruturada estritamente dentro das competências do Curso Técnico em Desenvolvimento de Sistemas.

\newpage

% ------------------------------------------------------------------------------
% 3. ÁREA DE ABRANGÊNCIA
% ------------------------------------------------------------------------------
\section{Área de Abrangência}

A área de abrangência do projeto compreende o ambiente operacional das feiras livres e feiras de produtores locais no contexto do Estado do Amazonas, com foco inicial de aplicação em feiras comunitárias e eventos de produtores sediados no município de Manaus.

O projeto não se restringe a um único espaço físico fixo, sendo concebido com arquitetura genérica para atender a diferentes feiras livres independentes. Os públicos e atores atendidos pelo sistema dividem-se em três perfis:

\begin{itemize}[leftmargin=1.5cm]
    \item \textbf{Produtores Locais / Feirantes:} Agricultores e feirantes cadastrados que disponibilizam seus catálogos semanais, recebem alertas de compras e realizam a separação das encomendas.
    \item \textbf{Clientes / Compradores:} Cidadãos que frequentam as feiras e utilizam o sistema para reservar antecipadamente produtos hortifrúti e artesanais.
    \item \textbf{Organização da Feira:} Administradores do espaço responsáveis por coordenar os pontos de entrega/retirada, ativar as edições da feira e acompanhar a logística geral.
\end{itemize}

O raio de atuação delimita-se à gestão de informações dentro do ciclo temporal que antecede a feira até o momento de encerramento das atividades do dia de atendimento presencial.

\newpage

% ------------------------------------------------------------------------------
% 4. OBJETIVOS E ESCOPO DO PROJETO
% ------------------------------------------------------------------------------
\section{Objetivos e Escopo do Projeto}

\subsection{Objetivo Geral}
Desenvolver uma aplicação web para a gestão de pedidos antecipados, catálogo de produtores locais e controle de pontos de retirada para feiras livres, estabelecendo um canal digital eficiente entre a agricultura familiar e os consumidores.

\subsection{Objetivos Específicos}
\begin{itemize}[leftmargin=1.5cm]
    \item Implementar o módulo de cadastro e gestão de produtores locais e seus respectivos catálogos de produtos;
    \item Desenvolver a vitrine pública de visualização de catálogo de produtos organizado por edição ou semana ativa da feira;
    \item Construir o mecanismo de montagem de pedidos on-line por meio de carrinho de compras sem intermediação de pagamentos eletrônicos;
    \item Disponibilizar o acompanhamento do status do pedido (\texttt{RECEBIDO}, \texttt{EM\_PREPARACAO} e \texttt{ENTREGUE}) com emissão de comprovante digital com código único;
    \item Gerar relatórios analíticos consolidados de vendas por produtor com isolamento estrito de dados.
\end{itemize}

\subsection{Delimitação do Escopo do MVP}
Para garantir a viabilidade técnica e a entrega funcional no prazo de 15 dias, os limites do sistema ficam formalmente estabelecidos:

\begin{itemize}[leftmargin=1.5cm]
    \item \textbf{Escopo Contemplado:} Cadastro de usuários (produtores, clientes e organização), catálogo semanal, carrinho de compras com persistência local, emissão de comprovante para retirada física, painel de controle do produtor com atualização de status, sistema interno de notificações e relatórios analíticos de vendas por feirante.
    \item \textbf{Não-Escopo:} O sistema não contempla integração com gateways de pagamento on-line via cartão ou PIX automático (o pagamento é realizado presencialmente na retirada da banca), nem rastreamento por GPS ou serviços de entrega terceirizada (delivery), operando estritamente com pontos físicos de coleta na feira.
\end{itemize}

\newpage

% ------------------------------------------------------------------------------
% 5. ESPECIFICAÇÃO DE REQUISITOS E REGRAS DE NEGÓCIO
% ------------------------------------------------------------------------------
\section{Especificação de Requisitos e Regras de Negócio}

\subsection{Requisitos Funcionais (RF)}

Os Requisitos Funcionais descrevem os comportamentos, fluxos operacionais e serviços que o sistema SISFEIRA disponibiliza aos seus usuários:

\vspace{0.3cm}
\noindent
\begin{spacing}{1.15}
\begin{tabularx}{\textwidth}{|c|p{4.2cm}|X|}
\hline
\textbf{Código} & \textbf{Identificador} & \textbf{Descrição do Requisito} \\ \hline
\textbf{RF01} & Cadastro de Produtores e Produtos & Permitir o cadastro de produtores e o gerenciamento autônomo dos produtos que compõem seu catálogo agrícola. \\ \hline
\textbf{RF02} & Cadastro de Clientes & Permitir o registro e a autenticação de clientes/compradores para realização de pedidos. \\ \hline
\textbf{RF03} & Catálogo Semanal da Feira & Disponibilizar a visualização pública dos produtos cadastrados para a edição ou semana ativa da feira. \\ \hline
\textbf{RF04} & Carrinho de Compras & Permitir que o cliente selecione múltiplos itens, defina quantidades e estruture seu pedido antes da confirmação. \\ \hline
\textbf{RF05} & Confirmação e Comprovante & Confirmar a solicitação de compra, gravando o pedido e gerando comprovante digital com código único para retirada. \\ \hline
\textbf{RF06} & Gestão do Status do Pedido & Permitir ao produtor acompanhar e atualizar o estado do pedido (recebido, em preparação e entregue). \\ \hline
\textbf{RF07} & Histórico de Pedidos & Disponibilizar ao cliente a consulta de compras anteriores e acompanhamento das solicitações em andamento. \\ \hline
\textbf{RF08} & Relatório de Vendas & Gerar relatórios analíticos consolidados de vendas e faturamento exclusivos para cada produtor. \\ \hline
\textbf{RF09} & Notificação de Novos Pedidos & Emitir alertas no painel do produtor e envio de notificações quando novos pedidos forem registrados. \\ \hline
\textbf{RF10} & Definição do Ponto de Retirada & Exibir e exigir a seleção do local físico de coleta do pedido previamente cadastrado para a feira. \\ \hline
\end{tabularx}
\end{spacing}

\newpage

\subsection{Requisitos Não Funcionais (RNF)}

Os Requisitos Não Funcionais estabelecem os critérios de qualidade técnica, restrições operacionais, segurança e usabilidade da aplicação:

\vspace{0.3cm}
\noindent
\begin{spacing}{1.15}
\begin{tabularx}{\textwidth}{|c|p{3.5cm}|X|}
\hline
\textbf{Código} & \textbf{Categoria} & \textbf{Descrição do Requisito} \\ \hline
\textbf{RNF01} & Desempenho & Garantir carregamento das páginas de catálogo em tempo inferior a 2 segundos em conexões padrão. \\ \hline
\textbf{RNF02} & Segurança & Proteger credenciais com hash criptográfico (\textit{bcrypt}) e blindar consultas SQL com parâmetros posicionais (\$1, \$2). \\ \hline
\textbf{RNF03} & Usabilidade & Fornecer fluxo de compra simplificado e intuitivo, realizável em poucos cliques e sem complexidade de navegação. \\ \hline
\textbf{RNF04} & Disponibilidade & Assegurar estabilidade e disponibilidade do servidor durante as janelas de funcionamento e picos da feira. \\ \hline
\textbf{RNF05} & Portabilidade & Estruturar layout responsivo otimizado para navegadores de smartphones (\textit{Mobile-First}). \\ \hline
\textbf{RNF06} & Escalabilidade & Suportar acessos concorrentes sem exaustão de conexões via \textit{Connection Pooling} no banco PostgreSQL. \\ \hline
\textbf{RNF07} & Manutenibilidade & Adotar arquitetura modular desacoplada em padrão REST, separando rotas, controladores, serviços e repositórios. \\ \hline
\textbf{RNF08} & Integridade Transacional & Garantir consistência estrita (ACID) na persistência de pedidos e reserva de itens por meio de transações atômicas. \\ \hline
\end{tabularx}
\end{spacing}

\newpage

\subsection{Regras de Negócio (RN)}

As Regras de Negócio definem as restrições e políticas de funcionamento que governam os processos da aplicação:

\vspace{0.3cm}
\noindent
\begin{spacing}{1.15}
\begin{tabularx}{\textwidth}{|c|p{3.8cm}|X|}
\hline
\textbf{Código} & \textbf{Identificador} & \textbf{Enunciado da Regra de Negócio} \\ \hline
\textbf{RN01} & Temporalidade do Catálogo & A disponibilização de produtos aos clientes é estritamente delimitada pela semana ou edição ativa da feira. \\ \hline
\textbf{RN02} & Ciclo de Vida dos Pedidos & O status do pedido evolui obrigatoriamente de forma sequencial e controlada: \texttt{RECEBIDO} $\rightarrow$ \texttt{EM\_PREPARACAO} $\rightarrow$ \texttt{ENTREGUE}. \\ \hline
\textbf{RN03} & Emissão de Comprovante & A confirmação do pedido pelo cliente gera obrigatoriamente um código identificador único para conferência na entrega. \\ \hline
\textbf{RN04} & Alerta Imediato ao Feirante & O registro de uma nova compra dispara automaticamente uma notificação persistida para cada produtor com itens no pedido. \\ \hline
\textbf{RN05} & Ponto de Atendimento & Nenhum pedido pode ser concluído sem a seleção prévia e obrigatória de um ponto físico de retirada homologado. \\ \hline
\textbf{RN06} & Isolamento de Produtores & O relatório analítico de vendas e faturamento deve consolidar exclusivamente as vendas pertencentes ao produtor logado. \\ \hline
\textbf{RN07} & Autonomia Administrativa & O cadastro e a precificação de produtos devem ser realizados diretamente pelo próprio feirante em sua interface gráfica. \\ \hline
\end{tabularx}
\end{spacing}

\newpage

% ------------------------------------------------------------------------------
% 6. MODELAGEM CONCEITUAL E DE DADOS
% ------------------------------------------------------------------------------
\section{Modelagem Conceitual e de Dados}

A modelagem do sistema traduz os requisitos de negócio em artefatos formais de engenharia de software, orientando o fluxo de dados e a persistência relacional.

\subsection{Diagrama de Casos de Uso}

O Diagrama de Casos de Uso explicita as fronteiras da aplicação, mapeando como os perfis de usuários (\textit{Cliente}, \textit{Produtor Local}, \textit{Organização da Feira} e o \textit{Serviço de Notificações}) interagem com as funcionalidades da plataforma.

\vspace{0.4cm}
\begin{center}
    \fbox{\parbox{0.9\textwidth}{\centering \vspace{2.5cm} \textbf{Figura 1 -- Diagrama de Casos de Uso do SISFEIRA} \\[0.2cm] \textit{(Inserir imagem gerada: diagrama\_casos\_de\_uso.png)} \vspace{2.5cm}}}
\end{center}
\vspace{0.4cm}

\subsection{Diagrama de Classes de Domínio}

O Diagrama de Classes representa as entidades fundamentais da lógica de negócio, seus atributos tipados e a multiplicidade das associações que estruturam a aplicação.

\vspace{0.4cm}
\begin{center}
    \fbox{\parbox{0.9\textwidth}{\centering \vspace{2.5cm} \textbf{Figura 2 -- Diagrama de Classes de Domínio} \\[0.2cm] \textit{(Inserir imagem gerada: diagrama\_classes.png)} \vspace{2.5cm}}}
\end{center}
\vspace{0.4cm}

\newpage

\subsection{Diagrama de Atividades (Ciclo de Vida do Pedido)}

O Diagrama de Atividades ilustra a esteira completa percorrida pelo pedido, desde a navegação na vitrine virtual pelo cliente até a finalização presencial na banca do feirante.

\vspace{0.4cm}
\begin{center}
    \fbox{\parbox{0.9\textwidth}{\centering \vspace{2.5cm} \textbf{Figura 3 -- Diagrama de Atividades e Ciclo de Vida do Pedido} \\[0.2cm] \textit{(Inserir imagem gerada: diagrama\_atividades.png)} \vspace{2.5cm}}}
\end{center}
\vspace{0.4cm}

\subsection{Dicionário de Dados (DD)}

O Dicionário de Dados formaliza o esquema físico implementado no PostgreSQL, especificando os tipos nativos, volumes estimados, rotinas de retenção e restrições de integridade.

\vspace{0.3cm}

% --- TABELA 1: USUÁRIOS ---
\begin{spacing}{1.0}
\small
\noindent
\begin{tabularx}{\textwidth}{|p{3.0cm}|X|}
\hline
\multicolumn{2}{|l|}{\textbf{ENTIDADE: Usuário}} \\ \hline
\textbf{Descrição:} & Cadastro de clientes, produtores e organizadores da feira. \\ \hline
\textbf{Nome da Tabela:} & \texttt{usuarios} \\ \hline
\textbf{Volume Esperado:} & $\approx$ 100 a 500 registros ativos na fase piloto. \\ \hline
\textbf{Tempo de Retenção:} & Permanente enquanto a conta mantiver-se ativa. \\ \hline
\textbf{Rotina de Limpeza:} & Inativação lógica conforme solicitação do titular (LGPD). \\ \hline
\end{tabularx}

\vspace{0.25cm}
\noindent
\begin{tabularx}{\textwidth}{|p{1.6cm}|p{1.9cm}|p{1.2cm}|p{0.7cm}|p{1.9cm}|X|}
\hline
\multicolumn{6}{|l|}{\textbf{ATRIBUTOS: \texttt{usuarios}}} \\ \hline
\textbf{Atributo} & \textbf{Nome Campo} & \textbf{Tipo} & \textbf{Tam.} & \textbf{Restrição} & \textbf{Descrição} \\ \hline
Identificador & \texttt{id} & UUID & 36 & PK, NOT NULL & Identificador único gerado por \texttt{gen\_random\_uuid()}. \\ \hline
Nome & \texttt{nome} & VARCHAR & 150 & NOT NULL & Nome civil completo ou razão social do produtor. \\ \hline
E-mail & \texttt{email} & VARCHAR & 150 & NOT NULL, UK & Login único utilizado para autenticação e emissão do JWT. \\ \hline
Hash Senha & \texttt{senha\_hash} & VARCHAR & 255 & NOT NULL & Hash seguro de autenticação gerado via bcrypt (RNF02). \\ \hline
Telefone & \texttt{telefone} & VARCHAR & 20 & NULL / CHECK & Contato com DDD; obrigatório para perfil PRODUTOR. \\ \hline
Perfil & \texttt{perfil} & VARCHAR & 20 & NOT NULL, CHECK & Papel de acesso: 'CLIENTE', 'PRODUTOR' ou 'ORGANIZACAO'. \\ \hline
Criação & \texttt{criado\_em} & TIMESTAMP & -- & DEFAULT NOW() & Data e hora de criação do registro no banco. \\ \hline
\end{tabularx}
\end{spacing}

\newpage

% --- TABELA 2: PRODUTOS ---
\begin{spacing}{1.0}
\small
\noindent
\begin{tabularx}{\textwidth}{|p{3.0cm}|X|}
\hline
\multicolumn{2}{|l|}{\textbf{ENTIDADE: Produto}} \\ \hline
\textbf{Descrição:} & Catálogo de itens hortifrúti e artesanais comercializados. \\ \hline
\textbf{Nome da Tabela:} & \texttt{produtos} \\ \hline
\textbf{Volume Esperado:} & $\approx$ 200 a 1.000 itens por edição de feira. \\ \hline
\textbf{Tempo de Retenção:} & Permanente para preservação do histórico de vendas. \\ \hline
\textbf{Rotina de Limpeza:} & Desativação lógica (\textit{Soft Delete}, \texttt{ativo = FALSE}). \\ \hline
\end{tabularx}

\vspace{0.25cm}
\noindent
\begin{tabularx}{\textwidth}{|p{1.6cm}|p{1.9cm}|p{1.2cm}|p{0.7cm}|p{1.9cm}|X|}
\hline
\multicolumn{6}{|l|}{\textbf{ATRIBUTOS: \texttt{produtos}}} \\ \hline
\textbf{Atributo} & \textbf{Nome Campo} & \textbf{Tipo} & \textbf{Tam.} & \textbf{Restrição} & \textbf{Descrição} \\ \hline
Identificador & \texttt{id} & UUID & 36 & PK, NOT NULL & Identificador universal único do produto. \\ \hline
Produtor & \texttt{produtor\_id} & UUID & 36 & FK, NOT NULL & Vínculo com \texttt{usuarios(id)} dono do item. \\ \hline
Nome & \texttt{nome} & VARCHAR & 100 & NOT NULL & Denominação comercial do item agrícola ou manufaturado. \\ \hline
Descrição & \texttt{descricao} & TEXT & -- & NULL & Detalhes sobre cultivo, apresentação e características. \\ \hline
Preço & \texttt{preco} & DECIMAL & 10,2 & NOT NULL, CHECK & Valor unitário oficial de venda ($\ge 0$). \\ \hline
Unidade & \texttt{unidade\_medida} & VARCHAR & 20 & NOT NULL & Unidade padrão de comercialização (ex: 'kg', 'maço'). \\ \hline
Ativo & \texttt{ativo} & BOOLEAN & 1 & DEFAULT TRUE & Flag que controla a visibilidade pública no catálogo. \\ \hline
\end{tabularx}
\end{spacing}

\vspace{0.6cm}

% --- TABELAS 3 E 4: EDIÇÕES E PONTOS DE RETIRADA ---
\begin{spacing}{1.0}
\small
\noindent
\begin{tabularx}{\textwidth}{|p{3.0cm}|X|}
\hline
\multicolumn{2}{|l|}{\textbf{ENTIDADE: Edição da Feira}} \\ \hline
\textbf{Descrição:} & Ciclos temporais de realização das feiras presenciais. \\ \hline
\textbf{Nome da Tabela:} & \texttt{edicoes\_feira} \\ \hline
\textbf{Volume Esperado:} & $\approx$ 50 registros anuais. \\ \hline
\textbf{Tempo de Retenção:} & 5 anos para histórico e auditoria. \\ \hline
\textbf{Rotina de Limpeza:} & Arquivamento anual de edições encerradas. \\ \hline
\end{tabularx}

\vspace{0.25cm}
\noindent
\begin{tabularx}{\textwidth}{|p{1.6cm}|p{1.9cm}|p{1.2cm}|p{0.7cm}|p{1.9cm}|X|}
\hline
\multicolumn{6}{|l|}{\textbf{ATRIBUTOS: \texttt{edicoes\_feira}}} \\ \hline
\textbf{Atributo} & \textbf{Nome Campo} & \textbf{Tipo} & \textbf{Tam.} & \textbf{Restrição} & \textbf{Descrição} \\ \hline
Identificador & \texttt{id} & UUID & 36 & PK, NOT NULL & Identificador universal da edição. \\ \hline
Nome Edição & \texttt{nome\_identificador} & VARCHAR & 100 & NOT NULL & Rótulo descritivo do evento (ex: 'Semana 42 - 2026'). \\ \hline
Ativa & \texttt{ativa} & BOOLEAN & 1 & DEFAULT FALSE & Flag que define se a feira está aberta para pedidos (RN01). \\ \hline
\end{tabularx}

\vspace{0.6cm}

\noindent
\begin{tabularx}{\textwidth}{|p{3.0cm}|X|}
\hline
\multicolumn{2}{|l|}{\textbf{ENTIDADE: Ponto de Retirada}} \\ \hline
\textbf{Descrição:} & Locais físicos para recolhimento presencial das encomendas. \\ \hline
\textbf{Nome da Tabela:} & \texttt{pontos\_retirada} \\ \hline
\textbf{Volume Esperado:} & $\approx$ 5 a 20 pontos de logística cadastrados. \\ \hline
\textbf{Tempo de Retenção:} & Permanente. \\ \hline
\textbf{Rotina de Limpeza:} & Manutenção manual conforme layout da feira. \\ \hline
\end{tabularx}

\vspace{0.25cm}
\noindent
\begin{tabularx}{\textwidth}{|p{1.6cm}|p{1.9cm}|p{1.2cm}|p{0.7cm}|p{1.9cm}|X|}
\hline
\multicolumn{6}{|l|}{\textbf{ATRIBUTOS: \texttt{pontos\_retirada}}} \\ \hline
\textbf{Atributo} & \textbf{Nome Campo} & \textbf{Tipo} & \textbf{Tam.} & \textbf{Restrição} & \textbf{Descrição} \\ \hline
Identificador & \texttt{id} & UUID & 36 & PK, NOT NULL & Identificador universal do local de logística. \\ \hline
Nome & \texttt{nome} & VARCHAR & 150 & NOT NULL & Nome do ponto físico (ex: 'Tenda Principal'). \\ \hline
Endereço & \texttt{endereco} & TEXT & -- & NOT NULL & Logradouro completo com referências de retirada. \\ \hline
\end{tabularx}
\end{spacing}

\newpage

% --- TABELA 5: PEDIDOS ---
\begin{spacing}{1.0}
\small
\noindent
\begin{tabularx}{\textwidth}{|p{3.0cm}|X|}
\hline
\multicolumn{2}{|l|}{\textbf{ENTIDADE: Pedido}} \\ \hline
\textbf{Descrição:} & Intenções de compra confirmadas pelos clientes na feira. \\ \hline
\textbf{Nome da Tabela:} & \texttt{pedidos} \\ \hline
\textbf{Volume Esperado:} & $\approx$ 500 a 5.000 pedidos por mês. \\ \hline
\textbf{Tempo de Retenção:} & 2 anos para histórico e conciliação de faturamento. \\ \hline
\textbf{Rotina de Limpeza:} & Consolidação periódica em base histórica analítica. \\ \hline
\end{tabularx}

\vspace{0.25cm}
\noindent
\begin{tabularx}{\textwidth}{|p{1.6cm}|p{1.9cm}|p{1.2cm}|p{0.7cm}|p{1.9cm}|X|}
\hline
\multicolumn{6}{|l|}{\textbf{ATRIBUTOS: \texttt{pedidos}}} \\ \hline
\textbf{Atributo} & \textbf{Nome Campo} & \textbf{Tipo} & \textbf{Tam.} & \textbf{Restrição} & \textbf{Descrição} \\ \hline
Identificador & \texttt{id} & UUID & 36 & PK, NOT NULL & Identificador universal do pedido. \\ \hline
Cliente & \texttt{cliente\_id} & UUID & 36 & FK, NOT NULL & Vínculo com o comprador em \texttt{usuarios(id)}. \\ \hline
Edição & \texttt{edicao\_id} & UUID & 36 & FK, NOT NULL & Vínculo com a feira em \texttt{edicoes\_feira(id)}. \\ \hline
Ponto Retirada & \texttt{ponto\_retirada\_id} & UUID & 36 & FK, NOT NULL & Vínculo com o ponto em \texttt{pontos\_retirada(id)}. \\ \hline
Comprovante & \texttt{comprovante\_codigo} & VARCHAR & 20 & NOT NULL, UK & Código gerado para retirada física (RN03). \\ \hline
Status & \texttt{status} & VARCHAR & 30 & NOT NULL, CHECK & 'RECEBIDO', 'EM\_PREPARACAO' ou 'ENTREGUE' (RN02). \\ \hline
Valor Total & \texttt{valor\_total} & DECIMAL & 10,2 & NOT NULL & Montante financeiro total calculado no back-end. \\ \hline
Criação & \texttt{criado\_em} & TIMESTAMP & -- & DEFAULT NOW() & Data e hora do fechamento da solicitação. \\ \hline
\end{tabularx}
\end{spacing}

\vspace{0.6cm}

% --- TABELAS 6 E 7: ITENS DO PEDIDO E NOTIFICAÇÕES ---
\begin{spacing}{1.0}
\small
\noindent
\begin{tabularx}{\textwidth}{|p{3.0cm}|X|}
\hline
\multicolumn{2}{|l|}{\textbf{ENTIDADE: Item do Pedido}} \\ \hline
\textbf{Descrição:} & Detalhamento dos produtos e valores congelados na compra. \\ \hline
\textbf{Nome da Tabela:} & \texttt{itens\_pedido} \\ \hline
\textbf{Volume Esperado:} & $\approx$ 2.000 a 20.000 itens processados mensalmente. \\ \hline
\textbf{Tempo de Retenção:} & 2 anos, rigorosamente vinculado ao pedido pai. \\ \hline
\textbf{Rotina de Limpeza:} & Exclusão em cascata (\textit{ON DELETE CASCADE}). \\ \hline
\end{tabularx}

\vspace{0.25cm}
\noindent
\begin{tabularx}{\textwidth}{|p{1.6cm}|p{1.9cm}|p{1.2cm}|p{0.7cm}|p{1.9cm}|X|}
\hline
\multicolumn{6}{|l|}{\textbf{ATRIBUTOS: \texttt{itens\_pedido}}} \\ \hline
\textbf{Atributo} & \textbf{Nome Campo} & \textbf{Tipo} & \textbf{Tam.} & \textbf{Restrição} & \textbf{Descrição} \\ \hline
Identificador & \texttt{id} & UUID & 36 & PK, NOT NULL & Identificador universal do item do pedido. \\ \hline
Pedido & \texttt{pedido\_id} & UUID & 36 & FK, NOT NULL & Vínculo com o pedido pai em \texttt{pedidos(id)}. \\ \hline
Produto & \texttt{produto\_id} & UUID & 36 & FK, NOT NULL & Vínculo com o catálogo em \texttt{produtos(id)}. \\ \hline
Quantidade & \texttt{quantidade} & INTEGER & 4 & NOT NULL, CHECK & Volume demandado pelo comprador ($> 0$). \\ \hline
Preço Unitário & \texttt{preco\_unitario} & DECIMAL & 10,2 & NOT NULL & Preço unitário imutável congelado na transação. \\ \hline
\end{tabularx}

\vspace{0.6cm}

\noindent
\begin{tabularx}{\textwidth}{|p{3.0cm}|X|}
\hline
\multicolumn{2}{|l|}{\textbf{ENTIDADE: Notificação}} \\ \hline
\textbf{Descrição:} & Alertas internos aos produtores sobre novas encomendas. \\ \hline
\textbf{Nome da Tabela:} & \texttt{notificacoes} \\ \hline
\textbf{Volume Esperado:} & $\approx$ 1.000 a 5.000 alertas gerados por mês. \\ \hline
\textbf{Tempo de Retenção:} & 6 meses a partir da data de leitura. \\ \hline
\textbf{Rotina de Limpeza:} & Expulgo de notificações lidas com mais de 90 dias. \\ \hline
\end{tabularx}

\vspace{0.25cm}
\noindent
\begin{tabularx}{\textwidth}{|p{1.6cm}|p{1.9cm}|p{1.2cm}|p{0.7cm}|p{1.9cm}|X|}
\hline
\multicolumn{6}{|l|}{\textbf{ATRIBUTOS: \texttt{notificacoes}}} \\ \hline
\textbf{Atributo} & \textbf{Nome Campo} & \textbf{Tipo} & \textbf{Tam.} & \textbf{Restrição} & \textbf{Descrição} \\ \hline
Identificador & \texttt{id} & UUID & 36 & PK, NOT NULL & Identificador universal do alerta. \\ \hline
Produtor & \texttt{produtor\_id} & UUID & 36 & FK, NOT NULL & Vínculo com o feirante em \texttt{usuarios(id)} (RN04). \\ \hline
Pedido & \texttt{pedido\_id} & UUID & 36 & FK, NOT NULL & Vínculo com a encomenda em \texttt{pedidos(id)}. \\ \hline
Mensagem & \texttt{mensagem} & TEXT & -- & NOT NULL & Descrição resumida dos itens encomendados. \\ \hline
Lida & \texttt{lida} & BOOLEAN & 1 & DEFAULT FALSE & Flag de controle de pendências no painel. \\ \hline
\end{tabularx}
\end{spacing}

\newpage

\subsection{Índices de Otimização e Integridade Transacional}

Para assegurar o carregamento rápido do catálogo (RNF01) e acelerar as consultas agregadas de faturamento por produtor (RF08), foram implementados quatro índices explícitos na base relacional:

\begin{itemize}[leftmargin=1.5cm]
    \item \texttt{CREATE INDEX idx\_produtos\_produtor ON produtos(produtor\_id);}: Acelera a recuperação e gestão de catálogo no painel do feirante.
    \item \texttt{CREATE INDEX idx\_pedidos\_cliente ON pedidos(cliente\_id, criado\_em DESC);}: Otimiza a exibição do histórico cronológico de compras do consumidor.
    \item \texttt{CREATE INDEX idx\_notificacoes\_nao\_lidas ON notificacoes(produtor\_id) WHERE lida = FALSE;}: Índice parcial que agiliza o contador de alertas pendentes do painel.
    \item \texttt{CREATE INDEX idx\_itens\_pedido\_agrupamento ON itens\_pedido(produto\_id);}: Garante alta performance nas junções SQL durante o cálculo dos relatórios de vendas (RN06).
\end{itemize}

Além da indexação, a integridade do sistema em momentos de concorrência é garantida pelo encapsulamento do fechamento de pedidos em transações atômicas nativas (\texttt{BEGIN}, \texttt{COMMIT} e \texttt{ROLLBACK}) operando com clientes dedicados do \textit{Connection Pool}, assegurando a persistência simultânea e consistente do cabeçalho do pedido, itens associados e notificações geradas.

\newpage

% ------------------------------------------------------------------------------
% 7. METODOLOGIA E ATIVIDADES PRÁTICAS
% ------------------------------------------------------------------------------
\section{Metodologia e Atividades Práticas}

O desenvolvimento do projeto SISFEIRA apoiou-se em um ecossistema integrado de ferramentas modernas, leves e de código aberto, cobrindo o planejamento de tarefas, prototipação visual, modelagem declarativa, controle de versão e persistência de dados.

\subsection{Trello (Gestão Ágil e Quadro Kanban)}
O Trello foi empregado como ferramenta de gestão visual baseada no método Kanban, organizando as demandas da equipe ao longo do ciclo de 15 dias:
\begin{itemize}[leftmargin=1.5cm]
    \item \textbf{Mapeamento de Backlog:} Registro estruturado das funcionalidades centrais (FEAT-01 a FEAT-07);
    \item \textbf{Fluxo de Trabalho:} Segmentação em colunas (\textit{A Fazer}, \textit{Em Andamento}, \textit{Revisão/Testes} e \textit{Concluído}), garantindo previsibilidade diária;
    \item \textbf{Controle de Prazos:} Alinhamento rigoroso com os marcos intermediários estabelecidos pelo cronograma pedagógico.
\end{itemize}

\subsection{Figma (Prototipação de Interfaces e UI/UX)}
Plataforma vetorial utilizada para a concepção e validação dos protótipos de interface antes da escrita de código front-end:
\begin{itemize}[leftmargin=1.5cm]
    \item \textbf{Design Mobile-First:} Telas concebidas prioritariamente para visualização em smartphones de produtores e clientes;
    \item \textbf{Simulação de Fluxos:} Validação da esteira de navegação (catálogo, carrinho, checkout e emissão de comprovante);
    \item \textbf{Guia de Identidade Visual:} Padronização da paleta de cores (verde floresta e amarelo terroso) e tipografia institucional.
\end{itemize}

\subsection{Mermaid.js (Modelagem Conceitual e Diagramas como Código)}
Biblioteca baseada em JavaScript utilizada para geração automatizada de diagramas UML a partir de texto puro (\textit{Diagrams as Code}):
\begin{itemize}[leftmargin=1.5cm]
    \item \textbf{Diagramação Ágil:} Geração e manutenção de Casos de Uso, Classes de Domínio, Atividades e Modelo Entidade-Relacionamento;
    \item \textbf{Versionamento com o Código:} Integração dos diagramas diretamente ao repositório Git, evitando dependências de arquivos binários proprietários.
\end{itemize}

\subsection{Git e GitHub (Controle de Versão Distribuído)}
O Git foi utilizado localmente no terminal Linux para rastreabilidade de código, aliado ao GitHub para repositório remoto:
\begin{itemize}[leftmargin=1.5cm]
    \item \textbf{Commits Semânticos:} Registro atômico e explicativo de cada evolução do código-fonte;
    \item \textbf{Ramificação por Funcionalidade:} Trabalho isolado através de \textit{feature branches} (ex: \texttt{feature/FEAT-01-auth-users});
    \item \textbf{Segurança Operacional:} Capacidade de reversão ágil (\textit{rollback}) para pontos estáveis em caso de regressões.
\end{itemize}

\subsection{PostgreSQL e Utilitário psql (Banco de Dados Relacional)}
SGBD relacional corporativo executado nativamente em ambiente Debian Linux na porta padrão 5432:
\begin{itemize}[leftmargin=1.5cm]
    \item \textbf{Automação via Linha de Comando (\texttt{psql}):} Aplicação sequencial de migrações SQL (\textit{migrations}) e dados iniciais de homologação (\textit{seeds});
    \item \textbf{Conformidade ACID:} Garantia de consistência transacional e integridade referencial sem dependência de ORMs pesados.
\end{itemize}

\newpage

% ------------------------------------------------------------------------------
% 8. FASES DE EXECUÇÃO, METAS E CRONOGRAMA
% ------------------------------------------------------------------------------
\section{Fases de Execução, Metas e Cronograma}

\subsection{Metas do Projeto}
As metas representam as entregas intermediárias consolidadas no ciclo de 15 dias:
\begin{itemize}[leftmargin=1.5cm]
    \item \textbf{Meta 1 (Dias 1 a 3):} Conclusão do levantamento e detalhamento dos Requisitos (RF e RNF) e do esquema relacional de dados.
    \item \textbf{Meta 2 (Dias 4 a 6):} Elaboração e validação dos protótipos navegáveis das 7 telas centrais da aplicação no Figma.
    \item \textbf{Meta 3 (Dias 7 a 11):} Implementação dos módulos de back-end (API REST Node.js/Express com PostgreSQL) e interfaces em Vanilla JS.
    \item \textbf{Meta 4 (Dias 12 a 13):} Execução dos testes integrados simulando a jornada completa de compra, emissão de comprovante e relatórios.
    \item \textbf{Meta 5 (Dias 14 a 15):} Finalização do relatório técnico impresso e estruturação dos slides para apresentação oral.
\end{itemize}

\subsection{Fases de Execução}
A distribuição temporal do projeto ocorreu em 5 blocos sequenciais de 3 dias:

\vspace{0.3cm}
\noindent
\begin{spacing}{1.15}
\begin{tabular}{|p{8.5cm}|c|c|c|c|c|}
\hline
\textbf{Fases de Execução} & \textbf{D 1-3} & \textbf{D 4-6} & \textbf{D 7-9} & \textbf{D 10-12} & \textbf{D 13-15} \\ \hline
Fase 1: Levantamento de Requisitos e Modelagem DER & X &  &  &  &  \\ \hline
Fase 2: Prototipação de Interfaces e Planejamento &  & X &  &  &  \\ \hline
Fase 3: Implementação do Back-end e Banco de Dados &  &  & X &  &  \\ \hline
Fase 4: Implementação do Front-end e Integração &  &  &  & X &  \\ \hline
Fase 5: Testes Integrados, Documentação e Apresentação &  &  &  &  & X \\ \hline
\end{tabular}
\end{spacing}

\subsection{Cronograma de Atividades}
O cronograma operacional detalha as atividades desenvolvidas ao longo das duas semanas de trabalho:

\begin{table}[h!]
\centering
\begin{spacing}{1.15}
\begin{tabular}{|p{7.5cm}|c|c|c|c|c|}
\hline
\textbf{Atividades Previstas} & \textbf{D 1-3} & \textbf{D 4-6} & \textbf{D 7-9} & \textbf{D 10-12} & \textbf{D 13-15} \\ \hline
Levantamento de Requisitos e Elaboração do DER & X &  &  &  &  \\ \hline
Modelagem Entidade-Relacionamento do Banco & X &  &  &  &  \\ \hline
Criação dos Protótipos de Interface (Wireframes) &  & X &  &  &  \\ \hline
Configuração do Ambiente e Estruturação do Banco &  & X &  &  &  \\ \hline
Desenvolvimento das Rotas e API (Node.js/Express) &  &  & X &  &  \\ \hline
Desenvolvimento do Front-end (HTML5/CSS3/JS) &  &  &  & X &  \\ \hline
Integração Front-end / Back-end e Validação &  &  &  & X &  \\ \hline
Bateria de Testes Funcionais do MVP &  &  &  &  & X \\ \hline
Consolidação do Relatório Técnico e Apresentação &  &  &  &  & X \\ \hline
\end{tabular}
\end{spacing}
\end{table}

\newpage

% ------------------------------------------------------------------------------
% 9. ÓRGÃOS ENVOLVIDOS, MECANISMOS E NORMAS DE EXECUÇÃO
% ------------------------------------------------------------------------------
\section{Órgãos Envolvidos, Mecanismos e Normas de Execução}

\subsection{Órgãos Envolvidos}
A execução do presente projeto técnico envolveu as seguintes entidades:
\begin{itemize}[leftmargin=1.5cm]
    \item \textbf{Centro de Educação Tecnológica do Amazonas -- CETAM:} Entidade pública estadual responsável pelo suporte institucional, supervisão técnico-pedagógica e homologação avaliativa do componente curricular.
    \item \textbf{Equipe de Estudantes Desenvolvedores:} Formada por \varEstudanteUm, \varEstudanteDois\ e \varEstudanteTres, responsáveis pelo ciclo completo de engenharia (requisitos, modelagem, codificação, testes e documentação).
\end{itemize}

\subsection{Privacidade e Proteção de Dados (LGPD)}
Em conformidade com a Lei Geral de Proteção de Dados Pessoais (LGPD - Lei nº 13.709/2018), o sistema adota o princípio da coleta mínima (\textit{data minimization}):
\begin{itemize}[leftmargin=1.5cm]
    \item Coletam-se apenas os dados indispensáveis para o fluxo operacional: nome, e-mail para autenticação e telefone para contato via WhatsApp;
    \item Não há captura nem armazenamento de documentos de identificação pessoal (CPF, RG) ou registros de cartões de crédito;
    \item A desativação de cadastros de produtos opera via \textit{Soft Delete} preservando a integridade histórica dos pedidos sem expor dados desnecessários.
\end{itemize}

\subsection{Segurança, Simplificações do MVP e Análise de Riscos}
Para viabilizar a entrega de um protótipo ágil, simples e funcional em ambiente de desenvolvimento local (Debian Linux), foram adotadas concessões e proteções técnicas específicas:
\begin{itemize}[leftmargin=1.5cm]
    \item \textbf{Proteções Efetivas Implementadas:} Hashes seguros de senhas gerados com \textit{bcrypt}, prevenção contra injeção de SQL via consultas parametrizadas (\$1, \$2) no driver \texttt{pg}, sanitização nativa de templates no front-end (\texttt{escapeHTML}) para neutralizar ataques de XSS e isolamento de escopo por produtor (RBAC) validado estritamente no back-end.
    \item \textbf{Armazenamento de JWT em \texttt{localStorage} (Risco Aceito):} Facilita o controle de sessão no front-end sem ferramentas de compilação. Em caso de vulnerabilidade XSS no navegador, o token poderia ser lido por scripts maliciosos. Para um ambiente de produção real com dados financeiros, recomenda-se a migração obrigatória para cookies protegidos com as flags \texttt{HttpOnly}, \texttt{Secure} e \texttt{SameSite}.
    \item \textbf{Execução em Host Único Bare-Metal (Risco Aceito):} A API Node.js e o PostgreSQL executam no mesmo host Debian comunicando-se via loopback local (\texttt{127.0.0.1}). Essa decisão elimina a sobrecarga de orquestradores de contêineres, sendo adequada para a avaliação acadêmica, embora em ambiente de produção comercial demande segmentação de rede e certificados TLS/HTTPS.
\end{itemize}

\newpage

% ------------------------------------------------------------------------------
% 10. ORÇAMENTO
% ------------------------------------------------------------------------------
\section{Orçamento}

Pela utilização exclusiva de tecnologias abertas, software livre e execução em infraestrutura computacional pré-existente dos próprios estudantes, os custos diretos limitaram-se a despesas simbólicas de material de consumo para apresentação:

\begin{table}[h!]
\centering
\begin{spacing}{1.15}
\begin{tabular}{|c|p{7.5cm}|c|c|c|c|}
\hline
\textbf{Nº} & \textbf{Descrição} & \textbf{Unid.} & \textbf{Qtd.} & \textbf{V. Unit. (R\$)} & \textbf{Total (R\$)} \\ \hline
\multicolumn{6}{|l|}{\textbf{Material de Expediente e Impressão}} \\ \hline
1 & Impressão da Documentação Técnica Final & Folha & 30 & 0,50 & 15,00 \\ \hline
2 & Encadernação em Espiral com Capa Plástica & Unid. & 1 & 10,00 & 10,00 \\ \hline
3 & Papel Sulfite A4 para Rascunhos e Validação & Pacote & 1 & 10,00 & 10,00 \\ \hline
\multicolumn{5}{|r|}{\textbf{Subtotal Material:}} & 35,00 \\ \hline
\multicolumn{6}{|l|}{\textbf{Infraestrutura e Licenças de Software}} \\ \hline
1 & Licenças de Software (Node.js, PostgreSQL, VS Code) & -- & -- & 0,00 & 0,00 \\ \hline
2 & Ambiente de Hospedagem (Host Único Debian Linux) & -- & -- & 0,00 & 0,00 \\ \hline
\multicolumn{5}{|r|}{\textbf{Subtotal Software:}} & 0,00 \\ \hline
\multicolumn{5}{|r|}{\textbf{Valor Total do Projeto (R\$):}} & \textbf{35,00} \\ \hline
\end{tabular}
\end{spacing}
\end{table}

\newpage

% ------------------------------------------------------------------------------
% 11. REFERÊNCIAS
% ------------------------------------------------------------------------------
\section{Referências}

\noindent ASSOCIAÇÃO BRASILEIRA DE NORMAS TÉCNICAS. \textbf{NBR 6023}: Informação e documentação -- Referências -- Elaboração. Rio de Janeiro: ABNT, 2018.

\vspace{0.5cm}
\noindent BRASIL. \textbf{Lei nº 13.709, de 14 de agosto de 2018}. Lei Geral de Proteção de Dados Pessoais (LGPD). Brasília, DF: Presidência da República, 2018.

\vspace{0.5cm}
\noindent CENTRO DE EDUCAÇÃO TECNOLÓGICA DO AMAZONAS (CETAM). \textbf{Orientações Metodológicas para Elaboração de Projetos Técnicos}. Manaus: CETAM, 2026.

\vspace{0.5cm}
\noindent ELMASRI, Ramez; NAVATHE, Shamkant B. \textbf{Sistemas de Banco de Dados}. 7. ed. São Paulo: Pearson, 2018.

\vspace{0.5cm}
\noindent PRESSMAN, Roger S.; MAXIM, Bruce R. \textbf{Engenharia de Software}: uma abordagem profissional. 8. ed. Porto Alegre: AMGH, 2016.

\vspace{0.5cm}
\noindent POSTGRESQL GLOBAL DEVELOPMENT GROUP. \textbf{PostgreSQL 15 Documentation}. 2023. Disponível em: \textless https://www.postgresql.org/docs/ \textgreater. Acesso em: 15 fev. 2026.

\newpage

% ------------------------------------------------------------------------------
% APÊNDICE
% ------------------------------------------------------------------------------
\section*{Apêndice}
\addcontentsline{toc}{section}{Apêndice}

\noindent \textbf{Apêndice A -- Relação de Telas Mapeadas para Prototipação}

\begin{itemize}[leftmargin=1.5cm]
    \item \textbf{Tela 01:} Autenticação e Cadastro (Produtor e Cliente);
    \item \textbf{Tela 02:} Catálogo Público de Produtos da Edição da Feira;
    \item \textbf{Tela 03:} Carrinho de Compras e Seleção de Quantidades;
    \item \textbf{Tela 04:} Confirmação de Pedido e Emissão de Comprovante de Retirada;
    \item \textbf{Tela 05:} Painel Administrativo do Produtor (Gestão de Itens e Preços);
    \item \textbf{Tela 06:} Painel de Controle de Status dos Pedidos Recebidos;
    \item \textbf{Tela 07:} Relatório Analítico Consolidado de Vendas por Feirante.
\end{itemize}

\newpage

% ------------------------------------------------------------------------------
% ANEXO
% ------------------------------------------------------------------------------
\section*{Anexo}
\addcontentsline{toc}{section}{Anexo}

\noindent \textbf{Anexo A -- Termo de Referência do Sistema Web (Extrato Pedagógico)}

\vspace{0.5cm}
\noindent Documento de especificação curricular e diretrizes pedagógicas fornecido pela Coordenação do Curso Técnico em Desenvolvimento de Sistemas do CETAM para balizamento do projeto prático.

\end{document}